Appendix~\ref{sec:appendix_exp_affine} reports the exponential affine SDF specification and the corresponding simulation comparison.




\section{Exponential affine SDF specification and simulation comparison}
\label{sec:appendix_exp_affine}

This appendix reports the alternative functional-form specification and simulation comparison omitted from the compressed main text.
We use an exponential affine SDF model adapted from \cite{KN16}
\begin{equation}
	\label{eq:expaff_SDF}
	\Psi_{\tau,t}^{\mathrm{EA}} (\theta) = 
	\exp
	\left[
	-\sum_{h=\tau+1}^{t} 
	\left(
	\alpha + \log (1 + r_{f,h}) + \sum_{m \in M} \beta_{m} \cdot \log (1 + F_{m,h})
	\right)
	\right]
\end{equation}
with monthly (arithmetic) risk-free return $r_f$, monthly (arithmetic) zero-net-investment portfolio returns $F_m$, and parameter vector $\theta=(\alpha,\beta)$.

For the functional-form comparison, all other simulation ingredients follow Section~\ref{sec:simulation_study}; deal returns are generated by the simple linear or exponential affine SDF models described in Equations~\eqref{eq:linear_sdf} and~\eqref{eq:expaff_SDF}.
In the exponential affine case, we adjust the lognormally distributed error mean to zero by subtracting $0.5 \sigma^2$.

In Figure~\ref{fig:sim_exp_aff}, we compare the simple linear SDF model against the exponential affine specification.
For the single-factor model, the bias is minimized at the 12.5-year Horizon (150 months), where both functional forms perform almost identically and yield estimates extremely close to the true value of 1.0: the exponential affine yields a mean market beta of 0.999 (0.182), matching the simple linear's 0.994 (0.189).
For the two-factor (MKT+Alpha) model at this same 12.5-year Horizon, both functional forms again yield nearly identical parameter results, though with an attenuated market beta ($\hat{\beta}_{\mathrm{MKT}} \approx 0.85$) and a minor spurious positive alpha ($\hat{\alpha} \approx 0.001$).
In the two-factor case, the market beta bias is actually minimized later, around the 20-year Horizon ($\hat{\beta}_{\mathrm{MKT}} \approx 0.97$, $\hat{\alpha} \approx -0.0002$).

However, a salient difference emerges at extremely short compounding Horizons.
At a 1-month Horizon (closely resembling an NPV-only discounting framework), the exponential affine formulation deteriorates much more severely in both models, manifesting a larger upward bias (e.g., in the single-factor case, $\hat{\beta}_{\mathrm{MKT}} = 1.21$ vs.\ $1.13$) and elevated variance ($0.39$ vs.\ $0.27$) compared to the linear model.
Overall, the simulation evidence does not indicate a systematic finite-sample advantage of exponential affine SDFs \citep{GM07,BMP08} in our framework.
If anything, the simple linear model appears slightly better behaved in the reported designs.

\begin{figure}[H]
	\centering
	\includegraphics[width=0.99\textwidth]{simulation/bias_comparison5}
	\caption{\textbf{Simulation ($\beta_{\textrm{MKT}}$, $\alpha$):} Exponential affine vs.\ linear models.}
	\label{fig:sim_exp_aff}
\end{figure}
